\documentclass[10pt,a4paper]{amsart}
\usepackage[margin=19mm]{geometry}
\usepackage{amsmath,amssymb,mathtools}
\usepackage[scr=rsfs,cal=euler]{mathalfa}
\usepackage{xfrac}
\usepackage[unicode,hidelinks,bookmarks=false]{hyperref}
\newcommand{\ie}{i.e.\ }
\newcommand{\cf}{cf.\ }
\newcommand{\resp}{respectively\ }
\newcommand{\Subsection}{\subsection}
\providecommand{\nobreakdash}{\nobreak\hbox{-}}
\setlength{\emergencystretch}{2em}
\multlinegap0pt
\usepackage{bm}
\usepackage{bbm}
\usepackage{dsfont}
\usepackage{xspace}
\usepackage{cancel}
\usepackage{mathtools}
\usepackage{amsthm}
\makeatletter
\@ifundefined{theorem}{\newtheorem{theorem}{Theorem}}{}
\@ifundefined{lemma}{\newtheorem{lemma}[theorem]{Lemma}}{}
\makeatother
\newcommand{\Cau}{\mathcal{C}}
\newcommand{\C}{\mathbb{C}}
\newcommand{\Log}{\operatorname{Log}}
\newcommand{\N}{\mathbb{N}}
\newcommand{\Real}{\operatorname{Re}}
\newcommand{\Z}{\mathbb{Z}}
\title[Finite-sheeted Cauchy operator at rational corners]{Finite-sheeted Cauchy operator at rational corners}
\author{Louis Shuo Wang}
\date{Source: arXiv:2606.12722v1 (CC BY 4.0)}
\begin{document}
\maketitle

\noindent\emph{Source and licence.} Finite-sheeted Cauchy operator at rational corners. Version arXiv:2606.12722v1 (CC BY 4.0), distributed under the Creative Commons Attribution 4.0 International licence, as recorded in the arXiv licence metadata for this exact version and shown on the abstract page https://arxiv.org/abs/2606.12722v1. Full rights notice: licenses/arxiv-2606.12722-CC-BY-4.0.txt.\par\smallskip

\noindent\textit{Editorial scope.} Counted unit: the Lemma whose source label is \texttt{lem:polyhom-model} in the article (source0.tex lines 805--836, source label \texttt{lem:polyhom-model}), delivered together with the article's own complete original proof of it (source0.tex lines 838--887, one \texttt{proof} environment of 50 lines). No mathematics is written, repaired or re-derived: statement and proof are the article's own text line for line. Required context retained uncounted: lem:model (lines 762--776); eq:model (lines 765--770); eq:dalpha (lines 799--803). Every definition, display and label of those lines is retained unchanged. Omitted from this package and not redistributed: 1--761, 771--798, 804--804, 837--837, 888--2425. Nothing inside a retained excerpt is omitted or altered: every retained body is delivered exactly as the article's lines are written, so no omission receipt of any kind accompanies this package. Citations: the retained excerpts carry no citation command at all, so no bibliography entry of the article is transcribed and no reference is redistributed. Reference adaptation: none was needed and none was made; every reference inside the delivered unit resolves inside it, so no unresolved reference or citation is produced. Printed slips kept literal: no printed word, symbol, hypothesis or number of the counted statement or of its proof was altered. Layout and class: the article's own class is \texttt{sn-jnl} with packages including graphicx, multirow, amsmath, amssymb, amsfonts, amsthm, mathrsfs, appendix, xcolor, textcomp, manyfoot, booktabs, algorithm, algorithmicx; the wrapper is the lane's pinned \texttt{amsart} editorial wrapper at 10pt on A4, which restates the theorem-like environments the retained text needs and adds the article's own macro definitions that the retained text needs (\texttt{\textbackslash C}, \texttt{\textbackslash Cau}, \texttt{\textbackslash Log}, \texttt{\textbackslash N}, \texttt{\textbackslash Real}, \texttt{\textbackslash Z}), with extra wrapper packages (bm, bbm, dsfont, xspace, cancel, mathtools). The delivered numbering is the unit's own. Assets: no retained span contains a figure environment, a graphics-inclusion command, a PDF-page inclusion, a table or a TikZ picture; no image file is copied into this package, the package declares no asset dependency and no image file is redistributed with it. Licence: version arXiv:2606.12722v1 of this article is distributed under the Creative Commons Attribution 4.0 International licence, as recorded in the arXiv licence metadata for this exact version and shown on the abstract page https://arxiv.org/abs/2606.12722v1. A full rights notice is delivered at licenses/arxiv-2606.12722-CC-BY-4.0.txt. Characters: every retained excerpt is pure ASCII, so the delivered engine, XeLaTeX with its default Latin Modern text font, reports no missing character.
\begin{lemma}[Model interval expansion]\label{lem:model}
Let $\alpha\in\C$, $\Real\alpha>-1$, $\alpha\notin\Z_{\ge0}$, $a>0$. Then, as
$\xi\to0$ off $[0,a]$,
\begin{equation}\label{eq:model}
  \Cau_{(0,a)}[s^{\alpha}](\xi)
  =m(\alpha)\,(-\xi)^{\alpha}
   +\frac{1}{2\pi i}\sum_{n\ge0}\frac{a^{\alpha-n}}{\alpha-n}\,\xi^{n},
   \qquad m(\alpha)=\frac{1}{2i\sin\pi(\alpha+1)},
\end{equation}
as an asymptotic expansion. For fixed $\xi\notin[0,a]$ the left side is
holomorphic in $\alpha$ on $\{\Real\alpha>-1\}$; the two singular contributions
at any $\alpha\to n_0\in\Z_{\ge0}$ have canceling simple poles, with combined
residue zero, and the value at $\alpha=n_0$ contains the term
$-\tfrac{1}{2\pi i}\xi^{n_0}\Log(-\xi)$.
\end{lemma}

\begin{equation}\label{eq:dalpha}
  s^{\alpha}(\log s)^{h}=\partial_\alpha^{h}s^{\alpha},
  \qquad
  (-\xi)^{\alpha}\bigl(\Log(-\xi)\bigr)^{h}=\partial_\alpha^{h}(-\xi)^{\alpha}.
\end{equation}

\begin{lemma}[Polyhomogeneous model expansion]\label{lem:polyhom-model}
Let $\alpha\in\C$, $\Real\alpha>-1$, $h\in\N_0$, $a>0$. For fixed
$\xi\notin[0,a]$,
\begin{equation}\label{eq:dalpha-transform}
  \Cau_{(0,a)}\bigl[s^{\alpha}(\log s)^{h}\bigr](\xi)
  =\partial_\alpha^{h}\,\Cau_{(0,a)}[s^{\alpha}](\xi),
\end{equation}
and the asymptotic expansion of \eqref{eq:model} may be differentiated termwise
in $\alpha$. Consequently:
\begin{enumerate}
\item \textup{(Nonresonant, $\alpha\notin\Z_{\ge0}$.)} As $\xi\to0$,
\begin{equation}\label{eq:polyhom-nonres}
  \Cau_{(0,a)}\bigl[s^{\alpha}(\log s)^{h}\bigr](\xi)
  =\sum_{i=0}^{h}\binom{h}{i}m^{(i)}(\alpha)\,(-\xi)^{\alpha}
     \bigl(\Log(-\xi)\bigr)^{h-i}
   +\sum_{n\ge0}c_{n,h}\,\xi^{n},
\end{equation}
with constant coefficients $c_{n,h}$ depending on $a,\alpha,n$ (no logarithm
occurs in the regular part when $\alpha\notin\Z_{\ge0}$). The top logarithmic
power of the singular part is $h$, with coefficient $m(\alpha)$.
\item \textup{(Resonant, $\alpha=n_0\in\Z_{\ge0}$.)} As $\xi\to0$,
\begin{equation}\label{eq:polyhom-res}
  \Cau_{(0,a)}\bigl[s^{n_0}(\log s)^{h}\bigr](\xi)
  =\frac{-1}{2\pi i\,(h+1)}\,\xi^{n_0}\bigl(\Log(-\xi)\bigr)^{h+1}
   +\sum_{0\le h'\le h}c_{h'}\,\xi^{n_0}\bigl(\Log(-\xi)\bigr)^{h'}
   +\sum_{n\ge0}\widetilde P_{n,h}(\log\xi)\,\xi^{n},
\end{equation}
with explicit constants $c_{h'}$ and polynomials $\widetilde P_{n,h}$ of degree
$\le h+1$. The top logarithmic power is $h+1$, raised by one relative to the
density.
\end{enumerate}
\end{lemma}

\begin{proof}
For fixed $\xi\notin[0,a]$ the integrand of
$\Cau_{(0,a)}[s^{\alpha}](\xi)=\frac{1}{2\pi i}\int_0^a s^{\alpha}(s-\xi)^{-1}ds$
is, by \eqref{eq:dalpha}, differentiated in $\alpha$ to give
$s^{\alpha}(\log s)^{h}(s-\xi)^{-1}$, which is absolutely integrable on $(0,a)$
for $\Real\alpha>-1$ and locally uniformly in $\alpha$; hence
\eqref{eq:dalpha-transform}. The expansion \eqref{eq:model} is the sum over poles
of the Mellin transform of $(s-\xi)^{-1}$ and is holomorphic in $\alpha$ on
$\{\Real\alpha>-1\}\setminus\Z_{\ge0}$, so its terms may be differentiated in
$\alpha$, yielding the asymptotic expansion of the left side of
\eqref{eq:dalpha-transform}.

\emph{Nonresonant.} Apply $\partial_\alpha^{h}$ to \eqref{eq:model}. By the
Leibniz rule and the second identity of \eqref{eq:dalpha},
$\partial_\alpha^{h}[m(\alpha)(-\xi)^{\alpha}]
=\sum_{i=0}^{h}\binom{h}{i}m^{(i)}(\alpha)(-\xi)^{\alpha}(\Log(-\xi))^{h-i}$,
the singular part of \eqref{eq:polyhom-nonres}. 
For the regular part,
$\partial_\alpha^{h}[a^{\alpha-n}(\alpha-n)^{-1}\xi^{n}]$ is, for
$\alpha\notin\Z_{\ge0}$, a constant (a polynomial of degree $\le h$ in $\ln a$)
times $\xi^{n}$, giving $c_{n,h}\xi^{n}$; no $\log\xi$ occurs in the regular
part for nonresonant $\alpha$.

\emph{Resonant.} Fix $n_0\in\Z_{\ge0}$ and set $\eta=\alpha-n_0$. Near $\eta=0$,
$m(\alpha)=\frac{\mu_{-1}}{\eta}+\mu_1\eta+\cdots$ (odd Laurent expansion, since
$1/\sin$ is odd about its pole) with $\mu_{-1}=\frac{(-1)^{n_0+1}}{2\pi i}$, and
$(-\xi)^{\alpha}=(-1)^{n_0}\xi^{n_0}e^{\eta L}$, $L:=\Log(-\xi)$. The $n=n_0$
regular term is $\frac{\xi^{n_0}}{2\pi i}\,\eta^{-1}e^{\eta\ln a}$. Their sum,
the only part singular as $\eta\to0$, is
\[
  G(\alpha,\xi)
  =\mu_{-1}(-1)^{n_0}\xi^{n_0}\,\frac{e^{\eta L}-e^{\eta\ln a}}{\eta}
   +\bigl(\text{terms regular in }\eta\bigr),
\]
using $\frac{\xi^{n_0}}{2\pi i}=-\mu_{-1}(-1)^{n_0}\xi^{n_0}$ to combine the two
$\eta^{-1}$ residues (which therefore cancel, confirming
Lemma~\ref{lem:model}). 
Now
$\frac{e^{\eta L}-e^{\eta\ln a}}{\eta}=\sum_{q\ge0}\frac{\eta^{q}}{(q+1)!}
\bigl(L^{q+1}-(\ln a)^{q+1}\bigr)$, so
$\partial_\alpha^{h}G|_{\alpha=n_0}=h!\,[\eta^{h}]G$ contributes
$\mu_{-1}(-1)^{n_0}\xi^{n_0}\frac{1}{h+1}\bigl(L^{h+1}-(\ln a)^{h+1}\bigr)$. The
non-analytic leading term is
$\mu_{-1}(-1)^{n_0}\frac{1}{h+1}\xi^{n_0}L^{h+1}
=\frac{-1}{2\pi i(h+1)}\xi^{n_0}L^{h+1}$, since
$\mu_{-1}(-1)^{n_0}=\frac{(-1)^{2n_0+1}}{2\pi i}=\frac{-1}{2\pi i}$. The
regular-in-$\eta$ remainder of $G$ together with the other $n\neq n_0$ regular
terms supplies the lower powers $L^{h'}$, $h'\le h$, and the analytic series with
log-coefficients of degree $\le h+1$; this is \eqref{eq:polyhom-res}.
\end{proof}

\end{document}
